% !TEX root = ../G-design.tex
\section{はじめに}
\section{眼球運動計測実験（EyeLink）}

\subsection{概要}
% 300字程度。目的・手法・結果・主要な考察を要約する。
本実験では，2枚の顔画像から魅力的と感じる方を選択する選好判断課題を用い，
選択に至るまでの視線が意思決定にどのように関与するかを検討することを目的とした。
眼球運動計測装置EyeLink~IIを用いて課題遂行中の注視位置を記録し，
キー押しによる選択反応時刻を基準として，選択された顔を注視していた割合の時間変化を
MATLABにより解析した。その結果，選択直前に向けて選択画像への注視割合が上昇する傾向が観察された。
この結果は，視線が選好の形成過程を反映しているとする先行研究の知見と整合する。
% ↑ 300字程度になるよう、実際の結果に合わせて末尾を調整してください。

\subsection{目的}
% 300字程度以上。実験を行う理由・意義。仮説を明示する。
ヒトの意思決定は外界から得た感覚情報に基づいて形成される認知過程であり，
その過程を可視化することは行動理解において重要な課題である\cite{tohoku2023}。
とりわけ視線は，選好の形成と相互に影響し合うことが知られている。
Shimojoら\cite{shimojo2003}は，2つの顔から好ましい方を選ぶ課題において，
選択に至る直前に選択対象への注視が段階的に増加する「視線カスケード効果
（gaze cascade effect）」を報告した。この現象は，単に選好が視線を導くだけでなく，
注視そのものが選好の形成に寄与する可能性を示唆している。
そこで本実験では，選好判断課題における注視位置を計測し，
選択された画像への注視割合が選択反応時刻に近づくにつれて上昇するか否かを明らかにすることを目的とした。
仮説として，キー押し直前に向かって選択画像への注視割合が0.5を上回って増加すると予想した。

\subsection{方法}
% 400〜1000字程度以上。追試可能な水準で詳細に。
\subsubsection{実験参加者}
本実験は第2グループ（2班・11班）の学生から選ばれた20代男性１名を対象として実施した。

% 参加者の人数・年齢範囲・視力（矯正含む）などが記録されていれば追記してください。

\subsubsection{装置および刺激}
眼球運動の計測には眼球運動計測装置EyeLink~II（サンプリング周波数250~Hz）を用いた。
刺激としてディスプレイ上に2枚の顔画像を左右に並べて提示した。
被験者は各試行で魅力的と感じた方の顔をキー押しにより選択し，
その左右反応（左：52，右：54）と反応時間が記録された。
提示および反応取得の試行数は1名あたり20試行とした。

\begin{figure}[htbp]
  \centering
  \begin{subfigure}[b]{0.32\textwidth}
    \centering
    \includegraphics[width=\linewidth]{kigu.jpg}
    \caption{実験器具}
    \label{fig:eye_kigu}
  \end{subfigure}
  \hfill
  \begin{subfigure}[b]{0.32\textwidth}
    \centering
    \includegraphics[width=\linewidth]{jikkenfuukei.jpg}
    \caption{実験風景}
    \label{fig:eye_fuukei}
  \end{subfigure}
  \hfill
  \begin{subfigure}[b]{0.32\textwidth}
    \centering
    \includegraphics[width=\linewidth]{hikennsya.jpg}
    \caption{被験者の様子}
    \label{fig:eye_hikennsya}
  \end{subfigure}
  \caption{実験環境}
  \label{fig:eye}
\end{figure}
\clearpage
\subsubsection{手続き}
各試行では，画面中央に注視点を提示した後に2枚の顔画像を左右に提示し，
被験者は魅力的と判断した側の顔に対応するキーを押した。
選択反応（キー押し）の時刻を0~msとし，その1秒前（$-996$~msから0~msまで，4~ms間隔）の
区間を解析対象とした。

\subsubsection{データ解析}
解析にはMATLABを用いた。各試行のEyeLinkデータ（時間，測定眼情報，画面状態，
左目注視位置$x$座標等を含むCSVファイル）を\texttt{readmatrix}関数で読み込んだ。
計測できなかった時刻は欠損値として扱い，\texttt{ismember}関数を用いて
サンプリング時刻ごとにデータの有無を判定した。注視点提示時（画面状態1）の
$x$座標の平均を画面中心とみなし，画像提示時の$x$座標がこの中心より小さい場合を左（52），
大きい場合を右（54）と判定した。この注視方向がキー出力データの選択方向と一致した場合を1，
一致しなかった場合を0，欠損を$-1$として各試行の行列（\texttt{expos}）にまとめた。
続いて20試行分を統合し，各サンプリング時刻について欠損値を除外したうえで，
選択画像を注視していた割合（反応率）を算出した。
横軸を刺激提示（キー押し）からの時間，縦軸を選択画像に対する注視確率とするグラフを作成した。

% 追加課題（全班データの統合平均・標準誤差、検定）を行った場合はここに追記してください。

\subsection{結果}
% 結果と考察で800〜1000字程度以上。ここでは事実のみを記述する。
選択画像に対する注視確率の時間変化を図\ref{fig:eye_graph}に示す。

\begin{figure}[htbp]
  \centering
  \includegraphics[width=0.8\hsize]{graph.pdf}
  \caption{注視確率の時間変化}
  \label{fig:eye_graph}
\end{figure}
キー押しの約1秒前の時点では注視確率は約0.5であった。その後、キー押し前およそ800〜550msの区間では一時的に0.35前後まで低下したが、キー押し前約500ms以降は選択反応時刻に近づくにつれて単調に上昇し、キー押し直前には約0.75に達した。
\subsection{考察}
% 結果に基づき、根拠を示しながら推論する。事実と推測を分ける。
本実験では，選択反応時刻に近づくにつれて選択画像への注視確率が上昇する傾向が得られた。
これはShimojoら\cite{shimojo2003}が報告した視線カスケード効果と一致する結果であると考えられる。
選択直前に注視が特定の対象へ収束する現象は，視線が単に選好の結果を反映するだけでなく，
注視の継続自体が対象への選好を高める方向に働いている可能性を示唆する。
また，松宮ら\cite{tohoku2023}は意思決定の信号が眼球運動系に連続的に流れることを示しており，
本結果における注視確率の漸増も，意思決定過程が眼球運動に持続的に反映された結果として解釈できる。
% ↓ 推測であることを明示。
一方、キー押し前800〜550msの区間で注視確率が一時的に0.5を下回った点については、選択が確定する以前の段階では両画像を交互に比較する探索的な視線移動が生じ、選択される側とは逆の画像へ視線が向かう時間が相対的に多かった可能性が考えられる。その後の注視確率の上昇は、比較を経て選好が形成されるにつれ視線が選択対象へ収束していく過程を反映していると推測される。
なお，キー押し直後ではなく直前から注視確率が高まった点については，
選好が形成される前段階から視線の偏りが生じていた可能性が推測される。
本実験の限界として，試行数が20試行と限られ，被験者数も少ないため，
個人差や刺激画像の偏りの影響を排除しきれていない点が挙げられる。
今後は被験者数と試行数を増やし，刺激画像の魅力度を統制したうえで検証することが課題である。